% Generated draft: requires mathematical and source review.
\documentclass[11pt]{article}
\usepackage[margin=25mm]{geometry}
\usepackage{fontspec}
\setmainfont{Latin Modern Roman}
\usepackage{amsmath,amssymb,mathtools,unicode-math}
\setmathfont{Latin Modern Math}
\usepackage{enumitem,needspace,xurl,hyperref}
\setlength{\parindent}{0pt}
\setlength{\parskip}{5pt}
\setlength{\emergencystretch}{3em}
\newcommand{\N}{\mathbb{N}}
\newcommand{\Z}{\mathbb{Z}}
\newcommand{\Q}{\mathbb{Q}}
\newcommand{\R}{\mathbb{R}}
\newcommand{\C}{\mathbb{C}}
\newcommand{\F}{\mathbb{F}}
\newcommand{\A}{\mathbb{A}}
\newcommand{\PP}{\mathbb{P}}
\newcommand{\E}{\mathbb{E}}
\newcommand{\eps}{\varepsilon}
\newcommand{\dd}{\,\mathrm{d}}
\begin{document}
\section*{UnsolvedMath ID: 700651}
\textbf{Source problem code:} catalogue entry
\subsection{Definitions and mathematical statement}
For a positive integer \(r\), put \(A_r=\Z[x_1,\ldots,x_r]\), the commutative polynomial ring in \(r\) indeterminates, and define
\[
G_{n,r}=\operatorname{SL}_n(A_r)=\{M\in\operatorname{Mat}_{n\times n}(A_r):\det M=1\},\qquad n\geq 1.
\]
Under the standard homological convention, a group \(G\) has type \(\mathrm{FP}_m\), for an integer \(m\geq0\), if the trivial left \(\Z G\)-module \(\Z\) admits a projective resolution
\[
\cdots\longrightarrow P_2\longrightarrow P_1\longrightarrow P_0\longrightarrow\Z\longrightarrow0
\]
in which \(P_i\) is finitely generated for every \(0\leq i\leq m\); no finite-generation condition is imposed for \(i>m\).

\textbf{Question, subject to the convention qualification below.} Is it true that
\[
\forall r\geq1\ \forall m\geq0\ \exists N(r,m)\geq1\ \forall n\geq N(r,m),\qquad G_{n,r}\text{ has type }\mathrm{FP}_m?
\]
Here all quantified parameters are integers. This formulation reads ``\(n\) sufficiently large'' as allowing the threshold to depend on both \(r\) and \(m\). The source does not explicitly specify this dependence or whether \(r=0\), giving the ring \(\Z\), is included.

\textbf{Unresolved convention qualification.} The source defines \((\mathrm{FP})_m\) by reference to the definition of \((\mathrm{FP})_\infty\) in Kourovka problem 6.3, requiring finite generation only for resolution terms numbered at most \(m\). That referenced definition and its numbering are absent from the supplied record. Consequently, the standard definition above is a provisional interpretation, and its agreement with the intended source convention requires review. Under the displayed standard convention, every group has type \(\mathrm{FP}_0\); the source's low-degree assertions therefore do not establish the intended indexing.

The supplied source reports affirmative answers for \(m=0\), \(n\geq3\), attributed to A. A. Suslin, and for \(m=1\), \(n\geq5\), attributed to M. S. Tulenbayev and to U. Rehmann and C. Soul\'e. These assertions are retained using the source's indexing and have not been independently verified.
\subsection{Short English statement}
G. A. Noskov asks whether special linear groups over a polynomial ring over \(\Z\) satisfy each prescribed homological finiteness condition once the matrix size is sufficiently large. The source's resolution convention and the allowed dependence of the size threshold require clarification. The supplied registry labels the problem open; this is a supplied status, not an independently verified assessment.
\subsection{Sources}
\begin{itemize}
\item Supplied local UnsolvedMath registry record, permanent ID \(700651\), exact source code \texttt{catalogue entry}, title ``Kourovka Notebook Problem 10.36.'' Upstream revision: \texttt{d371ee917ac9647fbb06cd6fc529deec6f3dda82}. The record supplies the status ``open.''
\item G. A. Noskov, problem 10.36, \emph{The Kourovka Notebook (Unsolved Problems in Group Theory)}, 10th issue (1986); the supplied background also identifies the 21st issue edition (2026), arXiv:1401.0300v48. The supplied collection description names E. I. Khukhro and V. D. Mazurov as editors. The definition referenced in problem 6.3 is not included in the supplied evidence.
\item A. A. Suslin, \emph{Math. USSR Izvestiya} \textbf{11} (1977), 221--238. Supplied attribution for the affirmative answer with source indices \(m=0\), \(n\geq3\); no article title is supplied.
\item M. S. Tulenbayev, \emph{Math. USSR Sbornik} \textbf{45} (1983), 139--154. Supplied attribution for the affirmative answer with source indices \(m=1\), \(n\geq5\); no article title is supplied.
\item U. Rehmann and C. Soul\'e, in \emph{Algebraic \(K\)-theory, Proceedings of the Conference, Northwestern University, Evanston, Illinois, 1976}, Lecture Notes in Mathematics \textbf{551}, Springer, Berlin, 1976, 164--169. Also cited by the source for \(m=1\), \(n\geq5\); no contribution title is supplied.
\item These references and attributions are supplied evidence only; their contents and the current literature status have not been independently verified.
\end{itemize}
\end{document}
